% =====================================================================
%  Chapter 24 opener -- "Chapter at a Glance", three reading paths + the
%  ladder staircase, THE running example (the shared 15-node tree T, top
%  centroid 4), the who-owns-what comparison against Ch.12 (centroid find),
%  Ch.13 (=K DP), Ch.18 Fenwick, Ch.22 HLD and Ch.17 sack, the recognition
%  guide, a constraint->technique table, and notation.
%  Exports fig:c18-tree, fig:c18-taxonomy.
%  No \term{centroid} or \term{centroid decomposition} here -- both are
%  indexed by Ch.12 (\chref{c09}); \term the NEW object "centroid tree" only.
% =====================================================================

\begin{chaptertheorybox}{Chapter at a Glance}
Chapter~12 taught you to \emph{find} a centroid: the node whose removal leaves
no piece bigger than \(n/2\), located in one linear walk toward whatever side is
heavy (\(\thmref{c09-centfind}\)).  That chapter stopped on purpose at "one
centroid, and why it is safe to cut here."  This chapter cashes that promise.  Do
not stop at one cut: cut at the centroid, then \emph{recurse} on every leftover
piece, and you get a divide-and-conquer whose recursion tree has only
\(O(\log n)\) levels -- because each cut leaves pieces at most half the size.  That
recursion is \emph{centroid decomposition}, and the tree it produces (this
chapter's new object, the \term{centroid tree}) is the load-bearing structure of
the whole chapter.

\textbf{You will learn to}
\begin{itemize}[nosep]
  \item build the \term{centroid tree} of a tree in \(\bigO{n\log n}\) -- find the
        centroid of a component, mark it deleted, recurse on the surviving
        pieces; prove the height is \(\le\lceil\log_2 n\rceil\) because every level
        at least halves a component (\(\secref{c18-tree}\), solved CF~321C); \hfill\lv{3}--\lv{4}
  \item \emph{claim every path by its top centroid} and count pairs at distance
        \(\le K\) (or \(=K\)) with a per-subtree distance histogram and the
        "subtract the same-child pairs" correction (\(\secref{c18-count}\),
        solved CF~715C); \hfill\lv{4}--\lv{5}
  \item use one centroid as a \emph{distance table} for its whole component: touch
        a node in \(O(\log n)\) writes up its centroid ancestors, then answer
        nearest / farthest / aggregate queries by walking the same ancestors
        (\(\secref{c18-index}\), solved CF~293E); \hfill\lv{5}
  \item run divide-and-conquer on trees \emph{beyond} plain path sums -- merge
        "arms" at the centroid, prune by an invariant, and let each bucket carry a
        Fenwick or sorted array (back-citing Ch.~18 \chref{c14})
        (\(\secref{c18-beyond}\), solved CF~833D); \hfill\lv{5}
  \item read a specification and decide whether the job is centroid's at all, or
        belongs to HLD / DSU-on-tree / tree DP / Mo / link-cut
        (\(\secref{c18-boundary}\), the pick table). \hfill\lv{5}
\end{itemize}

\textbf{Prerequisites.} Chapter~12 (\chref{c09}) -- the centroid definition and the
linear find; this chapter treats \(\thmref{c09-centfind}\) as a black box and never
re-derives it.  Chapter~13 (\chref{c10}) -- the \(O(nK)\) subtree DP that counts
distance-\(K\) pairs, the rival this chapter beats to \(O(n\log n)\) for large
\(K\).  Chapter~16 (\chref{c11}) -- \(dist(u,v)\) by depth and LCA, used to turn a
path count into a distance count.  Chapter~18 (\chref{c14}) -- the Fenwick tree
that rides inside a centroid bucket (\(\secref{c18-beyond}\) and the Summit).
\end{chaptertheorybox}

% ---------------- Reading paths + ladder position ----------------
\begin{chaptertheorybox}{How to Read This Chapter (three paths)}
Level badges \(\lv3\)--\(\lv6\): \(\lv3\) is decomposing the running tree \(T\) by
hand, cutting at the centroid and recursing; \(\lv{4}\) is the \(O(n\log n)\) build
written with a correct deleted-marker discipline, or a distance-\(K\) histogram with
the same-subtree subtraction; \(\lv{5}\) is a per-centroid distance table driving an
online nearest/aggregate query, or an arm-merge with a pruning invariant; \(\lv{6}\)
is the Summit, a centroid decomposition with a Fenwick living inside every bucket.

\textbf{Build-and-count path (\(\lv3\)--\(\lv{5}\), "count over all paths").}
\secref{c18-tree} in full -- the recursion, the \(\bigO{\log n}\)-level proof, solved
CF~321C -- then \secref{c18-count}, where paths are claimed by their top centroid and
counted by distance histogram (CF~715C), capped by the Summit \(\probref{c18-summit}\).

\textbf{Distance-index path (\(\lv{4}\)--\(\lv{6}\), "one node touches \(\log n\) tables").}
\secref{c18-index} -- the per-centroid distance table, used offline to count pairs
within a distance bound (solved CF~293E: \(\le l\) edges and \(\le w\) total weight) and,
with updates, for the online nearest-marked pattern (the classic Xenia and Tree, whose
statement Ch.~12 previews, named only);
then \secref{c18-beyond}, where a bucket carries a Fenwick (Ch.~18 \chref{c14}) --
solved CF~833D, and the Summit.

\textbf{Boundary path (\(\lv{4}\)--\(\lv{5}\), "is this centroid's job?").}
\secref{c18-boundary}: the pick table that hands path \emph{range} queries to HLD
(\chref{c16}), subtree multiset unions to DSU-on-tree (\chref{c28}), per-node states
to tree DP (\chref{c10}), offline node-set ranges to Mo (\chref{c21}), and mutable
topology to link-cut (\chref{c23}).
\end{chaptertheorybox}

\begin{figure}[H]
\centering
\begin{tikzpicture}[diagram, x=8.5mm, y=6.5mm]
  \foreach \st/\lab in {1/lv1,2/lv2,3/lv3,4/lv4,5/lv5,6/lv6} {
    \fill[\lab!22, draw=\lab, line width=0.8pt] (\st-1,0) rectangle (\st,\st*0.52);
    \node[font=\scriptsize\bfseries, text=\lab] at (\st-0.5,\st*0.52-0.42) {L\st};}
  \node[annot, anchor=south west, font=\scriptsize]
    at (0.05,0.62) {prereq (Ch.12): find one centroid, linear walk};
  \node[annot, anchor=south west, font=\scriptsize]
    at (1.05,1.14) {24.1 cut at 4, recurse; O(n log n); parent array, depth log n};
  \node[annot, anchor=south west, font=\scriptsize]
    at (2.05,1.66) {24.1 solved CF 321C (output centroid tree); 24.2 distance histogram, subtract correction};
  \node[annot, anchor=south west, font=\scriptsize]
    at (3.05,2.18) {24.2 solved CF 715C; 24.3 per-centroid distance tables};
  \node[annot, anchor=south west, font=\scriptsize, text=lv5]
    at (4.05,2.70) {24.3 solved CF 293E; 24.4 colour/arm aggregates, solved CF 833D};
  \node[annot, anchor=south west, font=\scriptsize, text=lv6]
    at (5.05,3.22) {24.5 the pick table -- is it centroid's job; Summit CF 757G};
\end{tikzpicture}
\caption*{Chapter~24 on the six-step ladder: decomposing \(T\) by hand is
\(\lv3\); a correct build and a distance-\(K\) count are \(\lv{4}\); per-centroid
distance tables and an arm-merge with a Fenwick bucket are \(\lv{5}\); the Summit
(centroid decomposition with a Fenwick inside every centroid) is step~6.}
\end{figure}

% ---------------- The running example ----------------
\subsection*{The Running Example: One Tree, Decomposed Over and Over}

Every section in this chapter runs on the same unweighted tree \(T\) on
\(n=15\) nodes, fixed once here and authoritative (if a later trace disagrees,
this figure wins).  Edges, written as pairs:
\[
  (4,1)\ (4,2)\ (4,3)\ (1,5)\ (1,6)\ (6,11)\ (2,7)\ (2,8)\ (8,12)\ (3,9)\ (3,10)\ (9,13)\ (10,14)\ (14,15).
\]
Its centroid is \(4\): removing it leaves three pieces \(\{1,5,6,11\}\),
\(\{2,7,8,12\}\), \(\{3,9,10,13,14,15\}\) of sizes \(4,4,6\), all \(\le 7\); every
other node leaves a piece of \(9\) or more, so \(4\) is the \emph{unique} top
centroid.  Recurse into those three pieces and you get the \term{centroid tree}
drawn in \(\secref{c18-tree}\).

\begin{figure}[H]
\centering
\begin{tikzpicture}[diagram, scale=0.9, transform shape,
   tnode/.style={circle, draw=primaryblue!70, fill=primaryblue!8, line width=0.7pt,
                 minimum size=6.5mm, font=\scriptsize\bfseries, inner sep=1pt},
   cen/.style={circle, draw=warningred, fill=warningred!18, line width=1.1pt,
               minimum size=7mm, font=\footnotesize\bfseries, inner sep=1pt}]
  \node[cen] (n4) at (0,0) {4};
  \node[tnode] (n1) at (-5,1.7) {1};
  \node[tnode] (n2) at (0,2.6) {2};
  \node[tnode] (n3) at (5,1.7) {3};
  \node[tnode] (n5) at (-6.6,3.2) {5};
  \node[tnode] (n6) at (-4.2,3.2) {6};
  \node[tnode] (n11) at (-4.2,4.7) {11};
  \node[tnode] (n7) at (-1.3,4.1) {7};
  \node[tnode] (n8) at (1.3,4.1) {8};
  \node[tnode] (n12) at (1.3,5.6) {12};
  \node[tnode] (n9) at (3.4,3.2) {9};
  \node[tnode] (n13) at (3.4,4.7) {13};
  \node[tnode] (n10) at (6.6,3.2) {10};
  \node[tnode] (n14) at (6.6,4.7) {14};
  \node[tnode] (n15) at (8.1,5.6) {15};
  \draw (n4)--(n1) (n4)--(n2) (n4)--(n3);
  \draw (n1)--(n5) (n1)--(n6) (n6)--(n11);
  \draw (n2)--(n7) (n2)--(n8) (n8)--(n12);
  \draw (n3)--(n9) (n9)--(n13);
  \draw (n3)--(n10) (n10)--(n14) (n14)--(n15);
  \draw[warningred, line width=1.3pt, densely dashed] (-7.2,2.35) -- (-2.0,2.35);
  \draw[warningred, line width=1.3pt, densely dashed] (-2.6,3.05) -- (2.6,3.05);
  \draw[warningred, line width=1.3pt, densely dashed] (2.0,2.35) -- (7.6,2.35);
\end{tikzpicture}
\caption{The running tree \(T\) (\(n=15\)) and its centroid \(4\) (red).  Cutting
at \(4\) -- the dashed cuts -- splits \(T\) into the three pieces of sizes
\(4,4,6\); each of size \(\le \lfloor 15/2\rfloor\).  Those three pieces are
recursed into independently in \(\secref{c18-tree}\), producing the centroid tree.
Left piece \(A=\{1,5,6,11\}\), middle \(B=\{2,7,8,12\}\), right
\(C=\{3,9,10,13,14,15\}\).}
\label{fig:c18-tree}
\end{figure}

% ---------------- Which job, cost chart ----------------
\subsection*{Which Job? (the question-shape chart)}

\begin{center}
\small
\renewcommand{\arraystretch}{1.25}
\setlength{\tabcolsep}{3pt}
\begin{tabular}{@{}p{0.30\linewidth} p{0.17\linewidth} p{0.29\linewidth} >{\raggedright\arraybackslash}p{0.14\linewidth}@{}}
\toprule
\textbf{Question} & \textbf{Technique} & \textbf{Why not a single DFS/DP} & \textbf{Where} \\
\midrule
count pairs of nodes at distance \(=K\)/\(\le K\) & top-centroid claim + histogram & a rooted DP is \(O(nK)\); centroid is \(O(n\log n)\) & \secref{c18-count} \\
for every node, nearest node with property \(P\) & per-centroid distance table & no single root holds all-pairs distances & \secref{c18-index} \\
online mark a node / query nearest marked & centroid tree index & the marked set changes; \(O(\log n)\) tables per touch & \secref{c18-index} \\
derive the centroid tree / rank every node by its centroid depth & build by cut-and-recurse & it is the decomposition itself, not one rooted DP & \secref{c18-tree} \\
sum of distances from \(v\) to a set / range & centroid bucket + Fenwick & the set is dynamic; \(\log n\) buckets touched per node & \secref{c18-beyond} \\
aggregate over \emph{every} simple path & \emph{divide and conquer on trees} & the path, not the node, is the unit of work & \secref{c18-count} \\
path \emph{range} update/query, fixed tree & \emph{go to HLD} & chains answer range-on-path, not all-pairs & \chref{c16} \\
\bottomrule
\end{tabular}
\end{center}

\begin{insightbox}[One promise decides everything: split where the tree balances]
A rooted DP fixes a root and combines children; it cannot see a path that wanders
up through the root and back down without counting it once per root, and it has no
handle on "the set of all paths."  Centroid decomposition changes the unit of work
from \emph{node} to \emph{path}: a path is charged to the highest centroid it
touches, and at that centroid the path has a fixed midpoint, so its length is
just \(depth_u + depth_v\) over two independent arms.  Because every cut leaves
pieces \(\le n/2\), the recursion is shallow (\(O(\log n)\)) yet each node
participates in only \(O(\log n)\) centroid buckets -- which is what turns an
\(O(n^2)\) all-pairs job into \(O(n\log^2 n)\), and lets a whole component be
\emph{indexed} from one node (\(\secref{c18-index}\)).
\end{insightbox}

% ---------------- Taxonomy chart ----------------
\begin{figure}[H]
\centering
\begin{tikzpicture}[diagram, scale=0.95, transform shape,
    mnode/.style={circle, draw=#1, line width=0.9pt, fill=#1!8, minimum size=8mm,
                  font=\scriptsize\bfseries, text=#1!60!black}]
  \node[mnode=primaryblue] (a) at (0,0)   {cut};
  \node[annot, font=\scriptsize, anchor=west, text width=9.6cm, align=left] at (1.4,0) {find a centroid (Ch.12's find), delete it, recurse on each piece; the parent array records the parent centroid, depth is log-bounded (24.1)};
  \node[mnode=successgreen] (b) at (0,-1) {claim};
  \node[annot, font=\scriptsize, anchor=west, text width=9.6cm, align=left] at (1.4,-1) {a path is charged to the highest centroid on it; there it is depth(u)+depth(v) across two arms -- count by distance histogram (24.2, CF 715C)};
  \node[mnode=difforange]   (c) at (0,-2) {index};
  \node[annot, font=\scriptsize, anchor=west, text width=9.6cm, align=left] at (1.4,-2) {each centroid keeps a distance table for its component; touch a node = O(log n) writes up its centroid ancestors, walk them to answer nearest/farthest (24.3, CF 293E)};
  \node[mnode=warningred]   (d) at (0,-3) {merge};
  \node[annot, font=\scriptsize, anchor=west, text width=9.6cm, align=left] at (1.4,-3) {beyond counting: merge arms at the centroid, prune by an invariant, and let a bucket hold a Fenwick or sorted array (Ch.18) (24.4, CF 833D)};
  \node[mnode=softpurple]   (e) at (0,-4) {pick};
  \node[annot, font=\scriptsize, anchor=west, text width=9.6cm, align=left] at (1.4,-4) {is it centroid's job? all-pairs/path-statistic yes; path ranges -> HLD; subtree unions -> sack; per-node states -> DP; offline sets -> Mo; dynamic edges -> link-cut (24.5)};
\end{tikzpicture}
\caption{The chapter's five-state taxonomy.  \emph{cut} is the decomposition
itself; \emph{claim} and \emph{index} are the two ways it pays off (count paths by
their top centroid; index a component by its centroid distance table); \emph{merge}
is the general divide-and-conquer-on-trees template with a data structure in the
bucket; \emph{pick} is the boundary choosing between centroid and its neighbours.
Each solved problem sits on one row: CF~321C on \emph{cut}, CF~715C on \emph{claim},
CF~293E on \emph{index}, CF~833D on \emph{merge}, the Summit CF~757G on
\emph{index}+\emph{merge}.}
\label{fig:c18-taxonomy}
\end{figure}

% ---------------- Who owns what ----------------
\subsection*{Who Owns What (so nothing is learned twice)}

\begin{center}
\small
\renewcommand{\arraystretch}{1.25}
\begin{tabular}{@{}p{0.30\linewidth} p{0.24\linewidth} p{0.40\linewidth}@{}}
\toprule
\textbf{Idea} & \textbf{Owner} & \textbf{This chapter's share} \\
\midrule
centroid definition + linear find & Ch.~12 (\chref{c09}) & treated as a black box; 24.1 begins at "delete it and recurse", never re-derives existence \\
count distance-\(K\) pairs by subtree DP & Ch.~13 (\chref{c10}) & the \(O(nK)\) rival (\(\probref{c10-harvest-solved}\)); 24.2 upgrades it to \(O(n\log n)\) and explains when the DP still wins \\
diameter / farthest by two DFS & Ch.~12 (\chref{c09}) & cited, never recomputed, where a centroid farthest-table needs a baseline \\
Fenwick / prefix sums & Ch.~18 (\chref{c14}) & 24.4 and the Summit place one BIT per centroid bucket and reuse it verbatim \\
path = \(O(\log n)\) chain ranges & Ch.~22 (\chref{c16}) & contrasted in 24.5: HLD answers path \emph{ranges}, centroid answers counts over \emph{all} paths \\
subtree multiset union (sack) & Ch.~17 (\chref{c28}) & forward-cited rival for "merge the children's sets"; centroid is for path/whole-component work \\
\(dist(u,v)\) by depth and LCA & Ch.~16 (\chref{c11}) & 24.2 uses it to turn "distance \(K\)" into "sum of two depths at the top centroid" \\
dynamic topology (link/cut) & Ch.~25 (\chref{c23}) & explicitly out; centroid is built once on a fixed tree \\
\bottomrule
\end{tabular}
\end{center}

% ---------------- Recognition guide ----------------
\begin{recognitionbox}[Pattern card -- "a claim about paths, or about distance from a set"]
\reckey{statements that ask for something over \emph{all pairs} or \emph{all simple
paths}: ``count pairs of nodes at distance \(K\) (or \(\le K\), or divisible by
\(m\)'', ``find the longest/shortest path with a property that a subtree DP cannot
localise'', ``for each node, the distance to the nearest node of colour/type \(C\)'',
``maintain a dynamic set of marked nodes and answer nearest-marked after each
update'', ``sum of distances from \(v\) to all nodes in a set/range''.  The tell is
that the unit of work is a \emph{path} (two endpoints) or a \emph{whole-component
distance}, not a single rooted subtree.}
\recuse{
  \item count / aggregate over all paths by a property additive along the path -- claim by top centroid (\(\secref{c18-count}\));
  \item nearest / farthest node from a moving set, online -- per-centroid distance index (\(\secref{c18-index}\));
  \item a path statistic a rooted DP cannot fold (a path-gcd count, a colour-balanced path count) -- arm-merge D\&C on trees (\(\secref{c18-beyond}\)).
}
\recnotuse{
  \item a range \emph{update and query along one path} on a fixed tree -- HLD (\chref{c16});
  \item a value that is a \emph{state per node} with no distance dimension -- tree DP / rerooting (\chref{c10}, \chref{c20});
  \item \emph{merge the children's subtrees' sets} bottom-up -- DSU-on-tree/sack (\chref{c28});
  \item many offline queries over a node set with free order -- Mo on trees (\chref{c21});
  \item the tree's \emph{edges change} (link/cut) -- link-cut / Euler-tour trees (\chref{c23}).
}
One-line test: \emph{the problem quantifies over paths or needs a component-wide
distance reachable from a single touched node -- then centroid decomposition, not a
rooted DP and not a path-chain engine.}
\end{recognitionbox}

% ---------------- Constraint -> technique ----------------
\subsection*{Reading the Constraints (which technique the limits want)}

\begin{center}
\small
\renewcommand{\arraystretch}{1.25}
\begin{tabular}{@{}p{0.34\linewidth} p{0.40\linewidth} p{0.20\linewidth}@{}}
\toprule
\textbf{Constraint signature} & \textbf{What it rules in/out} & \textbf{Technique} \\
\midrule
count pairs at distance \(=K\), \(n\le2\cdot10^5\), \(K\) large & the \(O(nK)\) subtree DP times out; the path has no fixed root & \secref{c18-count} \\
for each node nearest same-colour node & all-pairs is \(O(n^2)\); a centroid table answers in \(O(\log n)\) probes & \secref{c18-index} \\
online paint black / query nearest black, \(Q\le10^5\) & each touch writes \(O(\log n)\) buckets; query walks them & \secref{c18-index} \\
longest path with a global invariant (gcd \(>1\)) & not a composable fold; merge arms at each centroid and prune & \secref{c18-beyond} \\
sum of distances from \(v\) to a permutation index range, swaps allowed & dynamic set per bucket; a Fenwick over positions per centroid & \secref{c18-beyond} \\
\bottomrule
\end{tabular}
\end{center}

% ---------------- Notation ----------------
\subsection*{Notation Used Only Here}

The tree is \(T=(V,E)\), \(|V|=n\), unweighted unless a section says otherwise;
\(dist(u,v)\) is the number of edges (or the weighted length) on their path.  A
\emph{component} is a connected set of live (not-yet-deleted) nodes; \emph{size} is
its node count.  \(\mathit{cen}(S)\) is a centroid of component \(S\): a node whose
removal leaves every piece of size \(\le |S|/2\) (\(\defref{c09-centroid}\), Ch.~12).
\(\mathit{par}_c(u)\) is the \term{centroid tree} parent of \(u\): the centroid of
the component one level up that split \(u\) away (\(\secref{c18-tree}\)); \(\mathit{cd}(u)\)
is \(u\)'s depth in that tree, \(0\le\mathit{cd}\le\lceil\log_2 n\rceil\).  An
\emph{arm} of a centroid \(c\) is the set of nodes in one piece left after removing
\(c\), viewed by its distance from \(c\); a distance histogram \(H_c[d]\) counts
nodes at distance \(d\) from \(c\).  A node is \emph{deleted} / \emph{blocked} once
it is used as a centroid and is skipped in every deeper recursion.  The \emph{top
centroid of a path} is the first centroid (highest in the centroid tree) lying on
that path -- the node that "claims" it (\(\secref{c18-count}\)).  Every section runs
on \(T\) above; if a table here and a trace there disagree, this opener is
authoritative.
